\chapter{Governed Analytics and Validation Program}
\label{ch:validation-program}

The distinction between a useful analytics platform and a scientific
vendor comparison is the evidence design. A large collection of launch
records can support exploratory model checking without supporting the
stronger claim that one device disagrees with another. The repositories
associated with this review therefore use a two-release program with a
versioned data contract, explicit use eligibility, and fail-closed
statistics. The program is tracked by the public
\href{https://github.com/D-sorganization/Tools/issues/4583}{Tools program
epic}; the computational authority is UpstreamDrift, while the desktop and
web interfaces in Tools consume its versioned results rather than
reimplementing estimators.

\section{Qualified Corpus Snapshot}

The controlled authority snapshot contains 261,666 rows from 27 traceable
sources. Size alone is misleading. Only 13,855 rows contain finite,
non-imputed values for all five primary launch inputs (ball speed, launch
angle, launch direction, spin rate, and spin axis). Before the per-source
minimum-size gate, 13,647 rows also contain carry, 13,532 lateral outcome,
9,850 descent angle, and 9,283 flight time. The strict model cohort contains
13,843 rows from 15 source families because each predicted outcome has its own
availability rule. Eligibility is evaluated by source, metric, provenance,
and intended use; a blank field is represented as \emph{unavailable}, never
silently filled merely to increase sample size.

The largest source is a 236,590-row ShotLink-derived table. It has no
complete five-input launch rows and no observed carry values. Its rights and
transformation lineage also require separate approval. It is therefore
quarantined: the source is \textbf{not eligible for vendor-model training},
public redistribution, or vendor-performance claims. It may be retained
internally for provenance review and for analyses that are independently
shown to match its permitted use. This rule prevents a large but unsuitable
table from dominating the apparent evidence base.

At the present qualification gate, no vendor-labelled surrogate is currently eligible
for production training. The corpus contains 11,699 TrackMan-labelled
rows and 9,298 strict launch-input rows, but it lacks a trusted repeating
player or session grouping. Foresight has four rows (two strict); FlightScope
has 2,794 rows but no strict five-input rows. The historical TrackMan surrogate
used unique shot identifiers as partitions, which cannot protect against
player, session, or source leakage; it is retained for audit with status
\path{retired_non_group_safe} and cannot be loaded, queried, or retrained.
A future model that passes the group-safe gate would be described as
\emph{vendor-comparable on its qualified source}, not as a reproduction or
certification of proprietary vendor software.

\begin{keypoint}
Every displayed statistic carries its source identity, row-selection rule,
units, method version, and backing row references. An unavailable result is a
valid result: it records why the requested inference is not supported.
\end{keypoint}

\section{Release A: Qualified Analytics Platform}

Release A delivers the professional analysis surface on data that passes the
declared gate. Its scope includes unit-aware scatter and distribution views,
model residuals, arbitrary variable relationships, within-player and
between-player covariance, session trends, uncertainty, exportable plots,
and reproducible project files. Player-level analysis requires an explicit
player identity column. Identity is never inferred from session, club, row
order, filename, or directory layout; missing and whitespace-only identities
are excluded and reported.

The statistical backend resides in UpstreamDrift behind a versioned contract.
That contract returns estimates, uncertainty, warnings, source lineage, and
machine-readable unavailable states. Tools supplies PyQt and React/Vite
views with parity in variable choices, units, formulas, backing-data export,
persistence, and explanatory text. Interface parity means the two clients
consume the same result semantics; it does not justify copying the
statistical implementation into TypeScript.

Exploratory all-pairs scans report multiplicity warnings. Within-player
analysis separates raw pooled association, player-mean-centred association,
between-player association, and per-player estimates; these quantities answer
different questions and must not be substituted for one another. Correlation
and principal-component loadings are descriptive, not causal. Predictive
work uses grouped splits whenever a trusted player or session identity exists
and otherwise states that leakage-resistant grouping is unavailable.

Strokes gained is similarly gated. A result labelled \emph{strokes gained}
requires a cited baseline table, a declared lie and distance state, and a
reproducible mapping from shot outcome to expected strokes. A distance-only
or target-error score without that backing is labelled a \emph{performance
proxy}; it must not borrow the strokes-gained name.

\section{Qualified Model-Comparison Results}

The Release A campaign evaluated seven open flight-model configurations on the
13,843-row strict cohort, producing 96,901 predictions with no execution
failures. Its primary estimands are source-stratified error, bias, association,
and residual relationships. Pooled metrics are retained only as descriptive
summaries because source composition, device, venue, ball, and acquisition
conditions are confounded.

For the Blackmore TrackMan-labelled carry cohort ($n=8{,}860$), the best tested
configuration was the Waterloo/Penner implementation: carry RMSE
$12.256$\,m, mean error $-4.093$\,m, and strong linear association with the
reported outcome. These are agreement statistics against the source's reported
carry, not an absolute-accuracy result and not evidence that the proprietary
TrackMan equations have been identified. For context, the best tested carry
configuration on one Garmin R10 source ($n=1{,}541$) had RMSE $7.150$\,m and
bias $-4.722$\,m, while one Garmin R50 source ($n=255$) had RMSE
$10.238$\,m and bias $+6.554$\,m. The sign change is one reason a pooled vendor
ranking would be scientifically unsafe.

The content-addressed calculation record is the private campaign's
\path{results/v2/model/model_campaign_manifest.json}. It binds the cohort,
predictions, source metrics, pooled descriptive metrics, residual-correlation
table, configuration, and source qualification outputs by SHA-256. The
companion \path{multisource_v2_model_report.md} states row selection and
limitations in human-readable form. Restricted rows remain in the private
authority; the public paper reports only qualified aggregates.

The qualified snapshot is pinned to private authority commit and three
independently verifiable digests:
\begin{quote}
\small\ttfamily
commit: 996ce08ef6d097a91cf06a0b7432ae64e2a1b13b\\
qualification: 679856edb290111a0106c0bcbd2b673f581fd6b0026409c4d643565b330709a\\
model campaign: b70935fb4209b3f3aa18c3129ae948c4e20da74536af4d88433e5f0875960ca7\\
capability manifest digest: 1906d19fcace3284ae99d9dd8de213a0\allowbreak{}e2dabe8c062e4d63edc06c98f7eaf92e
\end{quote}
These identifiers bind the reported aggregates without exposing restricted
shot rows.

\section{Release B: Paired-Device Scientific Validation}

Release B begins only after a preregistered paired-device protocol has been
executed. The same physical shot must be observed simultaneously by each
device under comparison, with device model and firmware, ball and club,
indoor/outdoor condition, environmental state, alignment, player identity,
shot identity, and measured-versus-estimated status recorded. Randomized
blocks should span speed, loft, spin, launch direction, impact location, and
representative mishits. A traceable reference system is required where a
criterion claim is made.

The primary estimands are paired bias, limits of agreement, error as a
function of magnitude and condition, repeatability, missingness, and
device-by-condition interaction. Bland--Altman analysis and hierarchical
models are more informative here than correlation alone: two devices can be
highly correlated while disagreeing systematically. Sample-size targets and
exclusion rules must be fixed before inspecting device differences. Firmware
changes form new strata rather than being pooled invisibly.

The preregistered confirmatory design uses exactly three robot-centre,
outdoor club/speed cells with 84 accepted shot pairs per cell (252 total).
Each pair requires observations from two devices and exactly one independent
reference; calibration must pass, device rows must be unique within a pair,
serial identifiers must be SHA-256 hashes, and all members share the common
trigger timestamp. The protocol, capture-row JSON Schema, Python validator,
and synthetic fixture are complete, but the paired observations have not yet
been collected. Accordingly, protocol readiness is not reported as a
completed device-validation result.

Neural-network vendor emulators remain internal research artifacts. Their
training manifests must pin the eligible sources, feature and target units,
splits, random seeds, software version, and evaluation cohort. They are
released only with held-out and out-of-distribution performance, uncertainty
or applicability bounds, and a model card that forbids device-emulation or
certification claims. Until the paired-device gate is satisfied, Release B
comparisons remain unavailable rather than being approximated from unrelated
shots.

\section{Reproducibility and Claim Discipline}

The repository fleet records exact Git revisions for the statistical
authority, user interfaces, paper, public data facade, and private data
authority. A release manifest binds those revisions to the data manifest and
its checksums. Source files remain immutable; normalization and exclusions
are transformations with their own lineage. Project exports contain both the
request and complete backing rows so another analyst can reproduce the
displayed result without relying on GUI state.

These controls deliberately narrow the claims. Public launch-monitor records
can compare an open flight model with reported outcomes under stated source
conditions. They cannot, by themselves, identify a proprietary vendor model,
separate sensor error from aerodynamic-model error, or establish cross-device
agreement. Those questions belong to Release B and require paired evidence.
